% ============================================================================
%  SECCIÓN 2.1: ESTRUCTURA DE LOS MARCOS DEL PROYECTO
% ============================================================================

\section{Estructura de los marcos del proyecto}

Los marcos del proyecto constituyen una s\'intesis cr\'itica del estado del
conocimiento sobre el problema en estudio y aportan los elementos te\'oricos
necesarios y suficientes para explicar el objeto de estudio de la
investigaci\'on. Representan la descripci\'on, la explicaci\'on y el an\'alisis,
en un plano te\'orico, del problema central que trata la investigaci\'on, y
proporcionan los principios te\'oricos y los conceptos que sustentan el
proyecto \citep{arias2012, hernandez2014}.

La estructuraci\'on de los marcos del proyecto contempla cuatro marcos
interrelacionados, que en conjunto permiten fundamentar cient\'ificamente la
propuesta de soluci\'on. En la tabla se resumen los marcos que se desarrollan
en la siguiente secci\'on, de acuerdo con la gu\'ia de proyectos de grado de la
Universidad Privada Domingo Savio.

\begin{table}[H]
  \centering
  \caption{Estructura de los marcos del proyecto}
  \label{tab:estructura-marcos}
  \contenidoConFuente{%
    \begin{tablaAPA}
      \begin{tabular}{@{}p{3.2cm}p{12.6cm}@{}}
        \toprule
        \textbf{Marco} & \textbf{Descripci\'on} \\
        \midrule
        Marco te\'orico &
        Proceso en el que se integran teor\'ias, estudios, antecedentes y enfoques
        te\'oricos de validez para el correcto encuadre del tema de investigaci\'on. \\
        \addlinespace
        Marco conceptual &
        Asignaci\'on del significado preciso y contextualizado a los conceptos,
        t\'erminos y variables principales del tema de investigaci\'on. \\
        \addlinespace
        Marco hist\'orico &
        Recopilaci\'on de los antecedentes hist\'oricos referentes al tema de
        estudio, citados en orden cronol\'ogico. \\
        \addlinespace
        Marco referencial &
        Recopilaci\'on de las normas legales (leyes, decretos, resoluciones y
        reglamentos) relacionadas con el tema de investigaci\'on. \\
        \bottomrule
      \end{tabular}
    \end{tablaAPA}%
  }{Elaboraci\'on propia en base a la revisi\'on documental.}
\end{table}

El marco te\'orico, por su car\'acter integrador, agrupa a su vez los siguientes
elementos, que se desarrollan en la secci\'on siguiente:

\begin{itemize}
  \item Antecedentes de la investigaci\'on, que sit\'uan el punto de partida
    del estudio.
  \item Rese\~na de la instituci\'on, que describe el \'ambito acad\'emico en el
    que se desarrolla el proyecto.
  \item Bases te\'oricas, que integran las teor\'ias y los conceptos
    cient\'ificos que sustentan la soluci\'on propuesta.
  \item Sistema de variables, que relaciona los conceptos te\'oricos con los
    aspectos medibles del objeto de estudio.
  \item Definici\'on de t\'erminos b\'asicos, desarrollada en el marco
    conceptual.
\end{itemize}

La elaboraci\'on de los marcos se realiz\'o siguiendo el m\'etodo hist\'orico
l\'ogico y el an\'alisis y la s\'intesis, declarados en la secci\'on de m\'etodos
de investigaci\'on, mediante la revisi\'on documental de la literatura
especializada y la normativa vigente aplicable al objeto de estudio. El
desarrollo de cada marco se presenta a continuaci\'on, manteniendo una l\'ogica
deductiva que parte de los fundamentos generales hasta llegar a los elementos
particulares del sistema propuesto.
